% =====================================================================
% School Wits — ANSWER PAPER (QA) template
%
% Every QP file has exactly one QA partner. They are matched by the
% \examq header, which MUST be byte-identical in both files — same
% paper-id, same number, same topic string, same marks.
%
% Each question here contributes two things:
%   1. \begin{mstab}  — the official mark scheme grid
%   2. \begin{ansbox} — the worked solution / model answer
%
% See README.md in this folder for the JSON shape these map to.
% =====================================================================

\documentclass[11pt,a4paper]{article}
\usepackage{amsmath}
\usepackage{amssymb}
\usepackage[table]{xcolor}
\usepackage{tabularx}
\usepackage{multirow}
\usepackage{geometry}
\usepackage{enumitem}
\usepackage{parskip}
\usepackage{titlesec}
\usepackage{tocloft}
\usepackage[most]{tcolorbox}

\geometry{margin=2.5cm}

\titleformat{\section}{\large\bfseries}{}{0em}{}
\setlength{\cftbeforesecskip}{6pt}

% ===== Pastel colours, one per question =====
% A question whose colour is undefined is a LaTeX compile error, so this
% runs to q24 — Maths D 4024/22 has 24 questions. The palette repeats from
% q10; only adjacent questions need to look different.
\definecolor{q1color}{HTML}{D6EAF8} % pastel blue
\definecolor{q2color}{HTML}{D5F5E3} % pastel green
\definecolor{q3color}{HTML}{FDEBD0} % pastel orange
\definecolor{q4color}{HTML}{F9E79F} % pastel yellow
\definecolor{q5color}{HTML}{E8DAEF} % pastel purple
\definecolor{q6color}{HTML}{FADBD8} % pastel red/pink
\definecolor{q7color}{HTML}{D1F2EB} % pastel teal
\definecolor{q8color}{HTML}{FCF3CF} % pastel cream
\definecolor{q9color}{HTML}{EBDEF0} % pastel lavender
\definecolor{q10color}{HTML}{D6EAF8}
\definecolor{q11color}{HTML}{D5F5E3}
\definecolor{q12color}{HTML}{FDEBD0}
\definecolor{q13color}{HTML}{F9E79F}
\definecolor{q14color}{HTML}{E8DAEF}
\definecolor{q15color}{HTML}{FADBD8}
\definecolor{q16color}{HTML}{D1F2EB}
\definecolor{q17color}{HTML}{FCF3CF}
\definecolor{q18color}{HTML}{EBDEF0}
\definecolor{q19color}{HTML}{D6EAF8}
\definecolor{q20color}{HTML}{D5F5E3}
\definecolor{q21color}{HTML}{FDEBD0}
\definecolor{q22color}{HTML}{F9E79F}
\definecolor{q23color}{HTML}{E8DAEF}
\definecolor{q24color}{HTML}{FADBD8}

% ===== Subject shorthands used in ANSWER papers =====
% Add Maths mark schemes use these; question papers never do. Anything
% added here must ALSO be taught to KaTeX in js/app.js's KATEX_MACROS, or
% it renders as raw source in the browser.
\newcommand{\dd}{\mathrm{d}}
\newcommand{\dydx}{\dfrac{\mathrm{d}y}{\mathrm{d}x}}
\newcommand{\ncr}[2]{{}^{#1}\mathrm{C}_{#2}}
\newcommand{\npr}[2]{{}^{#1}\mathrm{P}_{#2}}
\newcommand{\cosec}{\operatorname{cosec}}

% ===== Mark-scheme table =====
%
% PICK ONE OF THE TWO DEFINITIONS BELOW AND DELETE THE OTHER.
% The choice decides how many "&" separated cells each row has, which is
% exactly what the JSON parser counts. Mixing them inside one file, or
% writing 4-cell rows under the 3-column definition, breaks extraction.
%
% Grouped Question labels use NEGATIVE \multirow placed on the LAST row
% of each group, so the label draws on top of the cell tint (this avoids
% a multirow/columncolor clipping bug — that is why the row order looks
% "upside down": blank label first, real label last).

% ---- OPTION A: 3 columns — Question | Answer | Marks -----------------
% Used by: Physics 5054/11, Physics 5054/21
\newenvironment{mstab}[1]{%
  {\color{#1!75!black}\nobreak Mark Scheme}\par\nopagebreak\smallskip
  \arrayrulecolor{#1!70!black}%
  \renewcommand{\arraystretch}{1.3}%
  \tabularx{\textwidth}{%
    |>{\columncolor{#1!12}\raggedright\arraybackslash}p{1.9cm}|%
     >{\columncolor{#1!12}}X|%
     >{\columncolor{#1!12}\centering\arraybackslash}p{1.1cm}|}%
  \hline
  \rowcolor{#1!45} Question & Answer & Marks \\ \hline
}{%
  \endtabularx
  \arrayrulecolor{black}%
  \par\medskip
}

% ---- OPTION B: 4 columns — Question | Answer | Marks | Partial Marks -
% Used by: Add Maths 4037/11, 4037/21, Maths D 4024/12, 4024/22
% The optional first argument renames the 4th column (e.g. "Guidance").
%
% \newenvironment{mstab}[2][Partial Marks]{%
%   {\color{#2!75!black}\nobreak Mark Scheme}\par\nopagebreak\smallskip
%   \arrayrulecolor{#2!70!black}%
%   \renewcommand{\arraystretch}{1.3}%
%   \tabularx{\textwidth}{%
%     |>{\columncolor{#2!12}\raggedright\arraybackslash}p{1.9cm}|%
%      >{\columncolor{#2!12}}X|%
%      >{\columncolor{#2!12}\centering\arraybackslash}p{1.1cm}|%
%      >{\columncolor{#2!12}}p{4.2cm}|}%
%   \hline
%   \rowcolor{#2!45} Question & Answer & Marks & #1 \\ \hline
% }{%
%   \endtabularx
%   \arrayrulecolor{black}%
%   \par\medskip
% }

% ===== Alternative-method banner =====
% \altrow{qNcolor}{Alternative} — a full-width heading inside an mstab
% introducing a second way of earning the SAME part's marks. It supplies
% its own row break, so it takes no trailing \\ in the source. Its rows are
% conventionally coded in parentheses, e.g. \textbf{(M1)}.
%
% Use the 4-column form below with Option B; change {4} to {3} for Option A.
\newcommand{\altrow}[2]{\multicolumn{4}{|l|}{\cellcolor{#1!30}\textbf{#2}} \\ \hline}

% ===== Worked-solution box =====
\newtcolorbox{ansbox}[1]{
  colback=#1!8,
  colbacktitle=#1!30,
  colframe=#1!70!black,
  coltitle=black,
  fonttitle=\normalfont,
  title=Worked Solution,
  boxrule=0.6pt,
  arc=2mm,
  left=4pt, right=4pt, top=4pt, bottom=4pt,
  breakable
}

\newcommand{\topiclabel}[1]{\hfill\normalfont\small\itshape\textcolor{black!55}{#1}}

% Per-question head — same command, same four arguments, same rendered
% line as the question paper. This is the join key between the two files.
\newcommand{\examq}[4]{\section[#1 - Q#2]{#1 - Q#2 - #3 - #4}}

\begin{document}

\begin{center}
  {\LARGE \textbf{Cambridge O Level Physics 5054/21}}\\[6pt]
  {\large May/June 2025 -- Mark Scheme \& Model Answers}
\end{center}

\bigskip

%%%%%%%%%%%%%%%%%%%
% Must match the QP's \examq for Q1 character for character.
\examq{5054/21/M/J/25}{1}{Motion or Kinematics \textperiodcentered\ Forces or Dynamics}{9}
%%%%%%%%%%%%%%%%%%%

\begin{mstab}{q1color}
% A part worth 2 marks = 2 rows sharing one label. The label sits on the
% LAST row via \multirow{-2}, and the rows above it start with an empty
% first cell. \cline{2-3} separates rows inside a group (use \cline{2-4}
% under the 4-column definition); \hline closes the group.
 & curve upwards (at start) with decreasing gradient labelled A & B1 \\ \cline{2-3}
\multirow{-2}{1.9cm}{1(a)} & horizontal line labelled B & B1 \\ \hline
% A single-mark part is one row with the label inline. \newline inside a
% cell separates alternative acceptable answers.
1(b) & gradient / slope decreases \newline or \newline increase in speed for same time less at later times & B1 \\ \hline
% Subparts use the (letter)(roman) form, no space: 1(c)(i)
1(c)(i) & air resistance / drag & B1 \\ \hline
 & air resistance / upwards force increases (with increased speed) & B1 \\ \cline{2-3}
\multirow{-2}{1.9cm}{1(c)(ii)} & air resistance / upwards force equals / balances weight & B1 \\ \hline
\end{mstab}

\begin{ansbox}{q1color}
% Part headings inside the box are plain text followed by \\[4pt].
% They are NOT \section or \item — the parser splits on the (a) / (c)(i)
% prefix at the start of a line.
(a) Speed--Time Graph\\[4pt]
Draw a curve that starts at the origin, rises steeply at first and then bends
over with a steadily decreasing gradient, finally levelling off into a
horizontal straight line.

\vspace{4pt}
(b) How the Graph Shows the Acceleration Decreasing\\[4pt]
The gradient of a speed--time graph equals the acceleration. As time increases
the curve becomes less steep, so the acceleration decreases.

\vspace{4pt}
(c)(i) The Other Vertical Force\\[4pt]
Air resistance (drag).

\vspace{4pt}
(c)(ii) Why the Resultant Vertical Force Becomes Zero\\[4pt]
As the skydiver speeds up, the air resistance acting upwards increases until it
balances her weight, leaving no resultant vertical force.
\[ R = \sqrt{400^{2} + 100^{2}} = \boxed{410\ \text{N}} \]
\end{ansbox}

%%%%%%%%%%%%%%%%%%%
% An MCQ answer: no mark scheme table, just the worked solution.
% (This is how Physics 5054/11 is written — 40 ansbox, 0 mstab.)
\examq{5054/11/M/J/25}{2}{Momentum (Dynamics)}{1}
%%%%%%%%%%%%%%%%%%%

\begin{ansbox}{q2color}
\textbf{Answer: A}\\[4pt]
Momentum is mass $\times$ velocity, so its unit is kg\,m\,/\,s.
\end{ansbox}

\end{document}
